\begin{casebox}
\textbf{知识点小案例：线程数等于核心数也可能跑不满。}
特例是 8 个 worker 处理 8 个 chunk，但其中一个 chunk 是热点 key 密集区，处理时间是其他 chunk 的 5 倍。即使线程数刚好等于核心数，7 个 worker 也会提前空闲，最后等待最慢任务。再叠加线程迁移，worker 的热数据可能离开原 cache 或 NUMA 节点。由此推出规律：调度分析不能只看线程数，要看 task 粒度、负载分布、worker state、迁移和等待点。工程报告至少输出 submit、start、finish、worker id、CPU id 和输入规模，才能解释空闲来自调度、数据倾斜还是同步。
\end{casebox}
